\chapter{Αναλυτική Δήλωση Χρήσης Εργαλείων Τεχνητής Νοημοσύνης}
\label{app:ai}

Το παρόν παράρτημα εξειδικεύει τη σύντομη δήλωση που περιλαμβάνεται στις εισαγωγικές σελίδες της εργασίας και περιγράφει αναλυτικά τη χρήση εργαλείων γενετικής τεχνητής νοημοσύνης (\eng{generative AI}) κατά την εκπόνησή της. Οι οδηγίες εκπόνησης διπλωματικών εργασιών του Τμήματος επιτρέπουν τη χρήση τέτοιων εργαλείων ως βοηθημάτων για την παραγωγή κειμένου και τη βελτίωση της σύνταξης. Θέτουν όμως δύο ρητούς όρους: ο φοιτητής ελέγχει προσεκτικά, λέξη προς λέξη, το παραγόμενο κείμενο και φέρει ο ίδιος την τελική ευθύνη για το περιεχόμενο. Η παρούσα εργασία εκπονήθηκε εντός αυτού του πλαισίου. Η δήλωση που ακολουθεί καταγράφει με ακρίβεια τόσο την έκταση της χρήσης όσο και τους ελέγχους που πραγματοποιήθηκαν.

\section{Εργαλεία που χρησιμοποιήθηκαν}
\label{sec:ai-tools}

Χρησιμοποιήθηκαν εργαλεία τεχνητής νοημοσύνης και πράκτορες Codex της OpenAI. Τα εργαλεία αυτά λειτούργησαν ως διαδραστικοί βοηθοί προγραμματισμού και σύνταξης κειμένου (\eng{coding and writing agents}). Εκτελούσαν εργασίες που όριζε ο συγγραφέας, με βάση γραπτές προδιαγραφές και κριτήρια αποδοχής, και κάθε παραδοτέο τους υποβαλλόταν στους ελέγχους που περιγράφονται στην Ενότητα~\ref{sec:ai-verification}.

\section{Έκταση της χρήσης ανά δραστηριότητα}
\label{sec:ai-scope}

Η συμβολή των εργαλείων υπήρξε ουσιαστική και δηλώνεται χωρίς υποτίμηση. Συγκεκριμένα, τα εργαλεία παρείχαν σημαντική υποβοήθηση στις ακόλουθες δραστηριότητες:
\begin{itemize}
  \item στην υλοποίηση τμημάτων του λογισμικού — επιλυτές (\eng{solvers}), συντεταγμένες, πίνακες αποκοπής (\eng{pruning tables}), αυτοματοποιημένες δοκιμές και εργαλεία ελέγχου — πάντοτε με βάση τις προδιαγραφές και υπό την καθοδήγηση του συγγραφέα·
  \item στην παραγωγή και αναθεώρηση προκαταρκτικών εκδοχών του ελληνικού κειμένου των κεφαλαίων, συμπεριλαμβανομένης της παρούσας δήλωσης·
  \item στην κατασκευή της αλυσίδας μετρήσεων και επαλήθευσης — σενάρια μετρήσεων, παραγωγή πινάκων από αποθηκευμένα αρχεία αποτελεσμάτων, εργαλεία ελέγχου ισχυρισμών — και στη διενέργεια εσωτερικών ελέγχων ποιότητας του κώδικα και του κειμένου.
\end{itemize}

Αντιθέτως, η ερευνητική κατεύθυνση και οι στόχοι, τα κριτήρια αποδοχής των αποτελεσμάτων και η τελική διατύπωση κάθε κεφαλαίου αποτελούν αποκλειστικά έργο του συγγραφέα. Το ίδιο ισχύει για τις αποφάσεις σχεδίασης και τεκμηρίωσης: την επιλογή των αλγορίθμων, την πολιτική κατά την οποία ένα αποτέλεσμα χαρακτηρίζεται \code{exact} (ακριβές) μόνο όταν προκύπτει από ολοκληρωμένη παραδεκτή αναζήτηση, καθώς και την απόφαση διατήρησης των τεκμηρίων πιστοποίησης που τεκμηριώνεται στο Κεφάλαιο~\ref{chap:validation}.

Τα εργαλεία δεν χρησιμοποιήθηκαν ως πηγή βιβλιογραφικών ή πειραματικών δεδομένων. Όλα τα αριθμητικά μεγέθη του κειμένου προέρχονται από αποθηκευμένα αρχεία αποτελεσμάτων, τα οποία παράγονται από καταγεγραμμένες, αναπαραγώγιμες εντολές. Οι βιβλιογραφικές αναφορές αντιστοιχούν σε πηγές που εντοπίστηκαν και επαληθεύτηκαν ανεξάρτητα από τα εργαλεία.

\section{Έλεγχοι που πραγματοποιήθηκαν}
\label{sec:ai-verification}

Το αντικείμενο της εργασίας περιλαμβάνει χαρακτηρισμούς υψηλής ευθύνης, όπως «βέλτιστος», «πλήρης» και «ακριβής», ενώ τα γλωσσικά μοντέλα μπορούν να προτείνουν διατυπώσεις γενικότερες ή εντυπωσιακότερες από όσες στηρίζουν τα τεκμήρια. Για τον λόγο αυτό, κάθε διατύπωση που προέκυψε με υποβοήθηση ελέγχθηκε έναντι των ετικετών κατάστασης των αποτελεσμάτων, των καταγεγραμμένων περιορισμών και των αποθηκευμένων τεκμηρίων. Ειδικότερα, πραγματοποιήθηκαν οι ακόλουθοι έλεγχοι:
\begin{itemize}
  \item Η πλήρης συλλογή αυτοματοποιημένων δοκιμών (\code{pytest}, 701 δοκιμές) εκτελέστηκε στην τελική κατάσταση του κώδικα και ολοκληρώθηκε χωρίς καμία αποτυχία.
  \item Κάθε αριθμητικός ισχυρισμός του κειμένου ελέγχθηκε μηχανικά έναντι των αποθηκευμένων αρχείων αποτελεσμάτων, μέσω του σεναρίου επαλήθευσης \code{verify_results.py} και του εργαλείου ελέγχου \code{thesis_audit.py}.
  \item Οι βαθιές γραμμές αποτελεσμάτων με κατάσταση \code{exact} διασταυρώθηκαν ανεξάρτητα με τον εξωτερικό επιλυτή Nissy 2.0.8, ο οποίος χρησιμοποιεί δικούς του πίνακες, όπως περιγράφεται στο Κεφάλαιο~\ref{chap:validation}.
  \item Ο συγγραφέας έλεγξε, διόρθωσε και τελικοποίησε το κείμενο λέξη προς λέξη και φέρει την πλήρη ευθύνη του περιεχομένου.
\end{itemize}

\section{Ευθύνη και διαφάνεια}
\label{sec:ai-responsibility}

Η χρήση υποβοηθητικών εργαλείων δεν μεταφέρει την ευθύνη ορθότητας στα εργαλεία. Η τελική ευθύνη για το περιεχόμενο, την ορθότητα των υλοποιήσεων, την εγκυρότητα των πειραματικών αποτελεσμάτων, την ακρίβεια των βιβλιογραφικών αναφορών και τη διατύπωση των συμπερασμάτων ανήκει αποκλειστικά στον συγγραφέα.

Για λόγους διαφάνειας, η έκταση της χρήσης είναι επαληθεύσιμη και όχι απλώς δηλούμενη. Το ιστορικό ανάπτυξης στο συνοδευτικό αποθετήριο κώδικα της εργασίας — τα μεταδεδομένα των υποβολών (\eng{commit metadata}) και τα αρχεία τεκμηρίωσης της ροής εργασίας — καταγράφει τα σημεία στα οποία συνεισέφεραν τα εργαλεία. Η παρούσα δήλωση συντάχθηκε ώστε να συμφωνεί με το ιστορικό αυτό.

Συνολικά, η χρήση των εργαλείων παρέμεινε εντός του πλαισίου που ορίζουν οι οδηγίες του Τμήματος: τα εργαλεία λειτούργησαν ως βοηθήματα για την παραγωγή και τη βελτίωση κώδικα και κειμένου, ενώ ο έλεγχος λέξη προς λέξη, η τελική διατύπωση και η πλήρης ευθύνη του περιεχομένου ανήκουν στον συγγραφέα.
